\subsection{Density of actual rigid points}\label{subsec:fc-line-rigid-density}

\begin{lemma}[Archimedean subgroup density]\label{lem:fc-archimedean-subgroup-density}
\lean{Subgroup.dense_of_no_min}
\mathlibok
A nontrivial multiplicative subgroup of the positive real numbers is dense
if its elements greater than $1$ have no least member.
\end{lemma}

\begin{lemma}[Density of the value range]\label{lem:fc-value-range-dense}
\lean{EnvelopeProperty.Berkovich.RelativeSpectrum.fc_denseRange_valuation}
\leanok
\uses{conv:fc-line, lem:fc-archimedean-subgroup-density}
Let $F$ be algebraically closed and let $v$ be nontrivial, in the equivalent
form that $v(t)>1$ for some $t\in F$. Then the range of $v$ is dense in
$\R_{\ge0}$.
\end{lemma}

\begin{proof}
The values of nonzero elements form a nontrivial multiplicative subgroup
of the positive reals. It has no least member greater than $1$: if such
a value $s>1$ were least, choose a square root of an element of value $s$.
Its value $u$ satisfies $u^2=s$, hence $1<u<s$, a contradiction.
Apply the Archimedean subgroup density lemma. Thus each interval with
positive endpoints contains a field value. An interval beginning at $0$ contains
an interval with positive endpoints, so the range is dense also in the
nonnegative real line.
\end{proof}

\begin{lemma}[Rigid approximations to Gauss points]\label{lem:fc-gauss-rigid-approximation}
\lean{EnvelopeProperty.Berkovich.RelativeSpectrum.fc_exists_rigid_tendsto_gauss, EnvelopeProperty.Berkovich.RelativeSpectrum.fc_gauss_mem_closure_rigid}
\leanok
\uses{lem:fc-value-range-dense, lem:fc-branch-end-convergence}
Every positive-radius Gauss point is the limit of actual evaluation points.
Every radius-zero Gauss point is itself an evaluation point.
\end{lemma}

\begin{proof}
For $r>0$, density of field values provides a strictly increasing sequence
$s_n$ of positive field values below $r$ with $s_n\to r$. Choose $c_n\in F$
with $v(c_n)=s_n$. The evaluation point at $a+c_n$ has value $s_n$ on
$T-a$. Its coordinate therefore approaches $r$ from below. Whole-branch
end convergence shows that these actual points tend to $\zeta_{a,r}$.
At radius zero the defining Gauss seminorm is evaluation at its centre.
\end{proof}

\begin{theorem}[Rigid-point density]\label{thm:fc-rigid-density}
\lean{EnvelopeProperty.Berkovich.RelativeSpectrum.fc_denseRange_rigid, EnvelopeProperty.Berkovich.RelativeSpectrum.fc_denseRange_rigid_norm}
\leanok
\uses{lem:fc-gauss-rigid-approximation, thm:fc-nested-disc-classification}
The actual evaluation points are dense in the analytic affine line.
\end{theorem}

\begin{proof}
Their closure contains all Gauss points by the preceding lemma. Every
analytic-line point is a Gauss point or a limit of actual Gauss points
by nested-disc classification. The closure of the evaluation points is
closed, so in either case it contains the point. It is therefore the
whole analytic line.
\end{proof}
